% ============================================================================
\section{Temple Simulado}\label{sec:temple}
% ============================================================================

\subsection{Idea}

El Temple Simulado (\emph{Simulated Annealing}, Kirkpatrick, Gelatt y Vecchi, 1983) se inspira
en el templado de materiales: si un vidrio o un metal fundido se enfría muy despacio, sus
partículas tienen tiempo de colocarse en la configuración de mínima energía; si se enfría de
golpe, quedan atrapadas en una configuración peor. En la analogía, la ``energía'' es el coste
de la solución y la ``temperatura'' $T$ es un número que controla cuánto desorden se tolera.

El algoritmo es una búsqueda local con una diferencia: cuando la vecina generada es
\emph{peor} que la solución actual no se rechaza siempre, sino que se acepta con probabilidad
\[
   P(\text{aceptar}) \;=\; e^{-\delta/T}, \qquad \delta = C(y)-C(x) \ \ge 0 ,
\]
donde $\delta$ es el empeoramiento. Esta fórmula tiene dos propiedades que son toda la gracia
del método:
\begin{itemize}
  \item cuanto mayor es el empeoramiento $\delta$, menor es la probabilidad de aceptarlo;
  \item cuanto mayor es la temperatura $T$, más fácil es aceptar empeoramientos.
\end{itemize}

\begin{center}
\small
\begin{tabular}{@{}lccc@{}}
\toprule
 & $\delta=0.1$ & $\delta=1$ & $\delta=5$ \\ \midrule
$T=5$ (caliente) & 0.98 & 0.82 & 0.37 \\
$T=1$ & 0.90 & 0.37 & 0.007 \\
$T=0.1$ (frío) & 0.37 & 0.00005 & $\approx 0$ \\
\bottomrule
\end{tabular}
\end{center}

Al principio $T$ es alta y el algoritmo se pasea casi libremente por el espacio de soluciones
(\textbf{diversifica}); a medida que $T$ baja, se va pareciendo a una búsqueda local que solo
acepta mejoras (\textbf{intensifica}). Una exposición breve y rigurosa del método puede
leerse en~\cite{bertsimas}; una más informal, en~\cite[cap.~2]{luke}.

\subsection{Esquema de los apuntes}

\begin{apuntes}[Esquema de los apuntes: Temple Simulado, algoritmo general]
\begin{tabbing}
\quad\=\quad\=\quad\=\quad\=\kill
Construir una solución inicial $x\in S$ y poner $x^\ast=x$\\
Determinar la temperatura inicial $T>0$\\
Repetir \{\\
\> $n=0$\\
\> Repetir \{\\
\>\> Generar una solución vecina $y\in V(x)$\\
\>\> Evaluar $y$ calculando $\delta=C(y)-C(x)$\\
\>\> Si ($\delta<0$) [$y$ es mejor que $x$]: \ $x=y$; \ si $C(x)<C(x^\ast)$ poner $x^\ast=x$\\
\>\> Si ($\delta\ge 0$) [$y$ no es mejor que $x$]:\\
\>\>\> Generar un número aleatorio $u\in(0,1)$\\
\>\>\> Si $u<\exp(-\delta/T)$ hacer $x=y$\\
\>\> $n=n+1$\\
\> \} Mientras $n\le N(T)$\\
\> Disminuir $T$\\
\} Mientras no se haya alcanzado el criterio de parada
\end{tabbing}
\end{apuntes}

Los apuntes señalan qué hay que decidir para aplicarlo. Por un lado la \emph{adaptación al
problema}: conjunto factible, cálculo del coste, estructura de entornos y solución inicial.
Por otro la \emph{estrategia de templado}: temperatura inicial, cómo disminuye, número de
iteraciones por temperatura y criterio de parada. Los apartados siguientes responden a cada
punto.

\subsection{Adaptación al problema de cubrimiento}

\paragraph{Conjunto factible y coste.} Las soluciones son siempre coberturas factibles y el
coste es $z$ (sección~\ref{sec:programa}). Como la solución lleva su coste actualizado,
$\delta$ es una simple resta.

\paragraph{Solución inicial.} La del constructivo voraz (criterio del cociente).

\paragraph{Vecindad.} El esquema pide ``generar una solución vecina'', es decir, \emph{una
vecina al azar}, no la mejor.

\begin{decision}[Decisión: la vecina se genera con la perturbación $\Om{k}$]
La vecina de $x$ se obtiene con la función \cod{perturbar} (sección~\ref{sec:perturbar}): se
quitan $k$ columnas al azar, se reparan las filas descubiertas con el voraz aleatorizado y se
eliminan redundantes. Razones:
\begin{itemize}[itemsep=0pt]
  \item La vecina es siempre factible, como exige nuestro diseño.
  \item Un movimiento más simple no sirve como vecindad \emph{aleatoria}: si se quita una
        columna al azar la solución casi siempre deja de ser factible, y si se añade una al
        azar la solución solo empeora. Un movimiento de inserción de $\D3$ elegido al azar
        casi siempre empeora mucho, porque la mayoría de las columnas no vuelve redundante a
        ninguna otra; elegirlos al azar desperdiciaría casi todas las iteraciones.
  \item Con $\Om{k}$ la diferencia de coste $\delta$ es pequeña y puede ser positiva o
        negativa, que es lo que el temple necesita para funcionar.
  \item El tamaño del cambio se regula con un solo número, $k$.
\end{itemize}
\end{decision}

\paragraph{Temperatura inicial.} La temperatura no tiene unidades ``naturales'': un valor
bueno depende de lo grandes que sean las diferencias de coste entre vecinas, que cambian de
una instancia a otra.

\begin{decision}[Decisión: fijar $T_0$ a partir de una probabilidad $p_0$]
En lugar de calibrar $T_0$ directamente, se calibra la probabilidad $p_0$ con la que se
quiere aceptar, al empezar, un empeoramiento ``típico''. El programa genera 100 vecinas de la
solución inicial, calcula el empeoramiento medio $\bar\delta^{+}$ de las que empeoran y
despeja $T_0$ de $e^{-\bar\delta^{+}/T_0}=p_0$:
\[
   T_0 \;=\; \frac{-\,\bar\delta^{+}}{\ln p_0}.
\]
Así el parámetro tiene un significado claro (``al principio acepto la mitad de los
empeoramientos'' si $p_0=0.5$) y la temperatura se adapta sola a cada instancia.
\end{decision}

\fragmento{temple.c}{67}{83}

\paragraph{Enfriamiento.} Se usa el enfriamiento \emph{geométrico}, $T\leftarrow\alpha T$ con
$0<\alpha<1$. Es el ejemplo que dan los apuntes~\cite{yanez} (``por ejemplo, $T=\alpha T$,
$0<\alpha<1$'') y el más habitual: baja deprisa al principio (cuando casi todo se acepta y
hay poco que ganar) y cada vez más despacio después.

\paragraph{Iteraciones por temperatura.} $N(T)=L$ constante: se generan $L$ vecinas con cada
valor de la temperatura.

\paragraph{Parada y recalentamiento.} El criterio de parada es el tiempo de CPU. Puede
ocurrir que la temperatura baje tanto que ya no se acepte ningún cambio mucho antes de que se
acabe el tiempo.

\begin{decision}[Decisión: recalentar si el sistema se congela]
Si durante 5 temperaturas seguidas no se acepta ningún cambio de coste, se considera que el
sistema está congelado: la temperatura vuelve a $T_0$ y la búsqueda se reanuda desde la mejor
solución conocida $x^\ast$. Sin esto, con un enfriamiento rápido el algoritmo pasaría casi
todo el tiempo parado en un óptimo local. Volver a $x^\ast$ (y no seguir desde la solución
actual) es una forma de intensificación.
\end{decision}

\subsection{Parámetros}

\begin{center}
\begin{tabular}{@{}llp{9.6cm}@{}}
\toprule
Nombre & Símbolo & Significado y efecto \\ \midrule
\cod{k} & $k$ & Columnas que se quitan para generar una vecina. Pequeño: vecinas muy
               parecidas, pasos cortos y baratos. Grande: pasos largos y más caros. \\
\cod{lrc} & --- & Tamaño de la lista de candidatos al reparar. Con 1 la reparación es
               determinista; con más, hay más variedad de vecinas. \\
\cod{p0} & $p_0$ & Probabilidad inicial de aceptar un empeoramiento típico. Fija $T_0$. Alta:
               empieza muy ``caliente'' (explora). Baja: empieza casi como una búsqueda local. \\
\cod{alfa} & $\alpha$ & Factor de enfriamiento. Cercano a 1: enfría despacio. \\
\cod{L} & $L$ & Vecinas generadas en cada temperatura. \\
\bottomrule
\end{tabular}
\end{center}

\begin{idea}[Cómo interactúan $\alpha$, $L$ y el tiempo]
Con $k=2$ el programa genera unas 84\,000 vecinas por segundo, es decir, unas 420\,000 en los
5 segundos. Eso son $420\,000/L$ temperaturas distintas, y la temperatura final es
$T_0\,\alpha^{420\,000/L}$. Por tanto $\alpha$ y $L$ no son independientes: lo que importa es
\emph{a qué temperatura se llega cuando se acaba el tiempo}. Si se llega demasiado caliente,
el algoritmo no ha llegado a intensificar; si se enfría demasiado pronto, se pasa el tiempo
recalentando. La calibración (sección~\ref{sec:cal-temple}) lo muestra con claridad.
\end{idea}

\subsection{Recorrido por el código}

El bucle principal es una transcripción del esquema de los apuntes. Los comentarios citan la
línea del esquema a la que corresponde cada instrucción.

\fragmento{temple.c}{87}{120}

Dos detalles de implementación:
\begin{itemize}
  \item ``$x=y$'' no se hace copiando la solución, sino \emph{intercambiando} las dos
        variables (línea \cod{tmp = x; x = y; y = tmp;}). Es instantáneo.
  \item Cuando la vecina tiene exactamente el mismo coste ($\delta=0$), $e^{0}=1$ y se acepta
        siempre: el algoritmo puede desplazarse por ``mesetas'' de soluciones de igual coste.
        Esos desplazamientos no se cuentan como cambios aceptados para decidir si el sistema
        está congelado.
\end{itemize}
